\documentclass{article}
\usepackage[UTF8]{ctex}
\usepackage{amsmath,amssymb,amsfonts}
\usepackage{enumitem}
\usepackage{geometry}
\geometry{a4paper, margin=1in}
% === 错题本排版优化（不参与题目标识） ===
\setlength{\parindent}{0pt}
\setlength{\parskip}{0.35em plus 0.10em minus 0.08em}
\setlength{\emergencystretch}{3em}
\setlist{itemsep=0.25em, topsep=0.40em, parsep=0pt}
\allowdisplaybreaks[2]
% === 错题本排版优化结束 ===
\begin{document}
\section*{武汉理工大学2024年高等代数真题}
\section*{126．武汉理工大学2024年高等代数试卷}
一．解答题．\\

\subsection{题目1}
1．（10 分）计算行列式

\[
D_{n}=\left|\begin{array}{ccccc}
a+1 & a+2 & a+3 & \cdots & a+n \\
a+2 & a+3 & a+4 & \cdots & a+1 \\
a+3 & a+4 & a+5 & \cdots & a+2 \\
\vdots & \vdots & \vdots & & \vdots \\
a+n & a+1 & a+2 & \cdots & a+n-1
\end{array}\right| .
\]

\subsection{题目2}
2．（15 分）已知非齐次线性方程组

\[
\left\{\begin{array}{l}
x_{1}+x_{2}+x_{3}+x_{4}=-1 \\
4 x_{1}+3 x_{2}+5 x_{3}-x_{4}=-1 \\
a x_{1}+x_{2}+3 x_{3}+b x_{4}=1
\end{array}\right.
\]

有三个线性无关的解．

\begin{itemize}
  \item[（1）] 求系数矩阵的秩．
  \item[（2）] 求 \(a, b\) 及解集合的秩与极大线性无关组。
\end{itemize}

\subsection{题目3}
3．（15 分）已知 \(A\) 为 4 阶矩阵，且有特征值 \(1,-1,2,-2\) ，求 \(\left|\frac{1}{16}\left(A^{*}\right)^{*}+A^{3} A^{*}-2 E\right|\) 。

\subsection{题目4}
4．（20 分）已知 \(W=\left\{A \mid \operatorname{tr}(A)=0, A \in \mathbb{R}^{2 \times 2}\right\}\) ．

\begin{itemize}
  \item[（1）] 证明 \(W\) 为 \(\mathbb{R}^{2 \times 2}\) 的子空间。
  \item[（2）] 求 \(W\) 的维数与一组基。
\end{itemize}

\subsection{题目5}
5．（20 分）已知数域 \(P\) 上的线性空间

\[
V=\left\{f(t) \mid f(t)=a_{0}+a_{1} t+a_{2} t^{2}, a_{0}, a_{1}, a_{2} \in P\right\} .
\]

\(\mathcal{T}\) 为 \(V\) 上的一个线性变换，且 \(\mathcal{T}(f(t))=f(2), f(t) \in V\) 。

\begin{itemize}
  \item[（1）] 求 \(\mathcal{T}\) 在基 \(1, t, t^{2}\) 下的矩阵 \(A\) 。
  \item[（2）] 求 \(\mathcal{T}\) 的特征子空间的维数与基。
\end{itemize}

\subsection{题目6}
6．（20 分）已知二次型 \(f\left(x_{1}, x_{2}, x_{3}\right)=x_{1}^{2}+4 x_{2}^{2}+4 x_{3}^{2}+2 t x_{1} x_{2}-2 x_{1} x_{3}+4 x_{2} x_{3}\) 为正定二次型．

\begin{itemize}
  \item[（1）] 求 \(t\) 的取值范围．
  \item[（2）] 求一个非退化线性替换化二次型为规范型。
\end{itemize}

\subsection{题目7}
7．（20 分）已知 \(A\) 为三阶实对称矩阵， \(1,-1,0\) 为 \(A\) 的特征值，\(\alpha_{1}=(1,-2,1)^{\mathrm{T}}\) 为 \(\lambda=1\) 对应的特征向量，\(\alpha_{2}=(-1, a, 1)^{\mathrm{T}}\) 为 \(\lambda=+1\) 对应的特征向量，求 \(a\) 的值以及矩阵 \(A\) 。
\end{document}
